The main objective of the internship was to work with
DreamDiffusion\cite{dreamdiffusion}, an EEG to image generator. In this part of
the internship it was asked to reproduce the DreamDiffusion results, and if
possible, to improve them, this would make enough material to make a paper.
However it showed many difficulties, the paper's results were very hard to
reproduce as the code was working in a not specified environment. The README
gave many steps in reproducing the paper but some were totally missing and the
GitHub issues are full of users complaining about it.

\subsubsection{The dataset}

DreamDiffusion used the CVPR40 dataset
\cite{CVPR40-1}\cite{CVPR40-2}\cite{CVPR40link}. \cref{fig:eeg-setup} (FROM
\cite{CVPR40-2}) shows the setup for recording the dataset. One person looks at
an image while its brain signals are being recorded. The cap used 128 electrodes
and the images were shown for 10s with a black screen of 10s between two images.
The images shown come from 40 of the classes of the ImageNet dataset.

To work with the dataset we don't use the raw recorded EEG, we use the
prefiltered EEGs with a bandpass from 9 to 95Hz.

\insertfigure
    {assets/images/cvpr40-eeg-setup}
    {width=0.8\textwidth}
    {Setup for CVPR40 dataset (IMAGE FROM \cite{CVPR40-2})}
    {fig:eeg-setup}

\subsubsection{The DreamDiffusion paper}

DreamDiffusion aims to be able to see what people see only from their brain
activity. This is still far from being a machine that can show what people can
dream of because the results are not that promising. The paper doesn't show any
metric, only examples of generated images and we discuss in the results section
about how far from perfect the model is.

The DreamDiffusion model is a built on top of pretrained Stable Diffusion
v1.5\cite{stable-diffusion}, but using a latent from a custom EEG autoencoder as
conditioning and adding a CLIP loss to align with. This EEG autoencoder is a
Masked Autoencoder (MAE) and is pretrained in a separated first time.

Ideas in this paper mostly come from the MinD-Vis \cite{mindvis} paper, tackling
the same problem of image generation but with fMRI instead of EEG as input.
DreamDiffusion took the code MinD-Vis code and changed it to take EEGs as input.
However, when MinD-Vis calculated their Fréchet Inception Distance \cite{fid}
(FID), a metric that assesses the image generation quality, DreamDiffusion
didn't, but in the MinD-Vis paper they precise that this score has poor
coherence in the context of a relatively small dataset, this is the case of
DreamDiffusion.

The metric DreamDiffusion used is the accuracy score of a pretrained ImageNet
classifier, it gives what part of the generated images were well classified in
the same class as the ground truth image really is. It gives a score of 45
percent to DreamDiffusion.

\cref{fig:dreamdiffusion-diagram} (FROM \cite{dreamdiffusion}) shows the whole
pipeline that is explained in the following.

\insertfigure
    {assets/images/dreamdiffusion-diagram}
    {width=\textwidth}
    {DreamDiffusion pipeline (IMAGE FROM \cite{dreamdiffusion})}
    {fig:dreamdiffusion-diagram}

\subsubsection{Training the EEG MAE}

A Masked Autoencoder is a model that was first used in \cite{mae}, the paper
shows the usage of these autoencoders with images, hiding patches of the image
during the training and teaching the model to reconstruct the base image. The
same idea is used here to have  

The MAE works as shown in \cref{fig:eeg-mae-diagram}. It is trained alone to
minimize the difference between the EEG sample and the reconstructed one and the
latent vector is supposed to carry all the information about the EEG sample but
in a different format. In the diagram $C$ is the number of channels of the EEG,
$T$ is the number of time steps in a sample, $d_e$ is the tokenizing output
dimension, $p$ is the number of patches the tokenizer makes, $r$ is the masking
ration ($0.5$ means that half the token will be removed), FC stands for fully
connected layer.

\begin{figure}[H]
    \centering
    \caption{EEG MAE diagram}
    \tikzset{
  >=Stealth,
  block/.style      = {draw, thick, shape=rectangle, minimum width=3.2cm,
                       minimum height=1.1cm, align=center, font=\footnotesize},
  group/.style      = {draw, dashed, rounded corners, inner sep=4pt},
  latent/.style     = {block, fill=gray!15},
  data/.style       = {block, fill=blue!10},
  model/.style      = {block, fill=green!10},
  loss/.style       = {block, fill=red!10},
  train/.style = {very thick},
}

\begin{tikzpicture}[node distance=1.6cm and 0.6cm]

% INPUT & PRE-PROC
\node[data] (sample) {EEG sample};
\draw ($(sample) + (-2.5, 0)$) edge[->, train] (sample);

% ENCODER
\node[model, below=of sample] (tokenize) {Tokenize};
\draw (sample) edge[->, train] node[midway, left] {$(C, T)$} (tokenize);
\node[model, below=of tokenize] (plus) {+};
\draw (tokenize) edge[->, train] node[midway, left] {$(d_e, \frac{T}{p})$} (plus);
\node[rotate=90, align=center, left=of tokenize] (pos_embed) {position\\embedding};
\path coordinate (midpoint0) at ($(pos_embed |- plus)$);
\draw (pos_embed) edge[train] (midpoint0)
      (midpoint0) edge[->, train] (plus);
\node[model, below= of plus] (masking) {Masking $(r)$};
\draw (plus) edge[->, train] node[midway, left] {$(d_e, \frac{T}{p})$} (masking);
\node[model, below= of masking] (vision) {Vision Transformer};
\draw (masking) edge[->, train] node[midway, left] {$(d_e, \frac{T\times r}{p})$} (vision);

% DECODER
\node[model, right=of vision] (fcdec) {FC};
\node[model, above=of fcdec] (unmasking) {Unmasking};
\draw (fcdec) edge[->, train] node[midway, right] {$(d_e', \frac{T\times r}{p})$} (unmasking);
\node[model, above= of unmasking] (visiondec) {Vision Transformer};
\draw (unmasking) edge[->, train] node[midway, right] {$(d_e', \frac{T}{p})$} (visiondec);
\node[model, above=of visiondec] (untokenize) {Untokenize};
\draw (visiondec) edge[->, train] node[midway, right] {$(d_e', \frac{T}{p})$} (untokenize);

% LATENT
\node[latent, below=of $(vision)!0.5!(fcdec)$] (latent) {latent};
\path coordinate (midpoint1) at ($(vision |- latent)$);
\draw (vision) edge[train] (midpoint1)
      (midpoint1) edge[->, train] (latent);
\path coordinate (midpoint2) at ($(fcdec |- latent)$);
\draw (latent) edge[train] (midpoint2)
      (midpoint2) edge[->, train] (fcdec);

\node[below= 0.5cm of latent] (mask) {mask};
\path coordinate (midpoint3) at ($(masking)+(-2,0)$);
\path coordinate (midpoint4) at ($(midpoint3 |- mask)$);
\draw (masking) edge[train] (midpoint3)
      (midpoint3) edge[train] (midpoint4)
      (midpoint4) edge[->, train] (mask);
\path coordinate (midpoint6) at ($(unmasking)+(2,0)$);
\path coordinate (midpoint5) at ($(midpoint6 |- mask)$);
\draw (mask) edge[train] (midpoint5)
      (midpoint5) edge[train] (midpoint6)
      (midpoint6) edge[->, train] (unmasking);

% OUTPUT
\node[data, above=of untokenize] (reconst) {Reconstructed};
\draw(untokenize) edge[->, train] node[midway, right] {$(C, T)$} (reconst);

% LOSS
\node[model, above=of sample] (mse) {Mean Square Error};
\draw (sample) edge[->, train] (mse)
      (reconst) edge[->, train] (mse);
\node[loss, right=of mse] (loss) {EEG Autoencoder loss};
\draw (mse) edge[->, train] (loss);


% GROUP BOXES (dashed) ---------------------------------------------------------
\node[group,fit=(tokenize)(plus)(masking)(vision)] (encgroup) {};
\node[rotate=90] at ($(encgroup.west) + (-0.8,0)$) {\small \textbf{EEG ENCODER}};
\node[group,fit=(fcdec)(unmasking)(visiondec)(untokenize)] (decgroup) {};
\node[rotate=-90] at ($(decgroup.east) + (0.8,0)$) {\small \textbf{EEG DECODER}};

\end{tikzpicture}
    \label{fig:eeg-mae-diagram}
\end{figure}

After a good amount of training with enough dropout we have a good EEG MAE, some
signals, being masked and reconstructed again are shown together in
\cref{fig:mae-reconstruct}. This is the reproduced result by Racim Menasria, an
older intern who worked on the subject with the same team.

\insertfigure
    {assets/images/reconstructed-eeg}
    {width=\textwidth}
    {Trained MAE results}
    {fig:mae-reconstruct}

\subsubsection{Training the diffusion model}

The second step was to train the stable diffusion model with this latent vector
as conditioning. In parallel we add a CLIP \cite{clip} loss to align the VAE
latent with the CLIP embeddings, that way, we can start from the pretrained
stable diffusion model as the conditionning from the pretrained checkpoint is
matching with CLIP embeddings.

The diffusion model comes from \cite{stable-diffusion}. It is a Variational
Autoencoder (VAE) followed by a diffuser, as shown in (FROM
\cite{stable-diffusion}) \cref{fig:stablediffusion-pipeline}.

\insertfigure
    {assets/images/stable-diffusion-pipeline}
    {width=\textwidth}
    {Stable Diffusion model (IMAGE FROM \cite{stable-diffusion})}
    {fig:stablediffusion-pipeline}

Contrary to a basic diffusion model\cite{diffusion}, this one is conditioned,
the attention layers in the U-NET take keys and values from the conditioning
latent vector.  This diffusion process have a loss for each step of the
diffusion, comparing the noised ground truth image with the denoised vector
$z_T$, plus, a loss from the CLIP embedding of the ground truth image.

To sum up, we can count three losses in this training that are combined

\begin{itemize}
    \item Diffusion loss: at each step of the diffusion, compare the denoised
    vector with the noised one

    \item VAE loss: compare $\mathcal{D}(\mathcal{E}(x))$ with $x$ for $x$ the
    ground truth image.
    
    \item CLIP alignement loss: align the CLIP image embedding of the ground
    truth image with the MAE latent vector. A projection layer transforms the
    EEG embedding from the pretrained MAE to an embedding that matches CLIP
    embedding of the ground truth image.
\end{itemize}

\subsubsection{The reproducibility}

Reproducing DreamDiffusion was a long task, not only the training took around 12
hours for the MAE but also around 20 hours for the diffusion model training, but
the initial GitHub repository of DreamDiffusion gave no checkpoint for the MAE
and no clear information on the training data for the MAE.

To train the MAE all that is needed is a lot of EEG signals, but as few EEG
datasets have 128 channels like the DreamDiffusion paper, data augmentation
techniques were used in the code. Concerning the datasets all the authors of
DreamDiffusion specified is that they took their data from the Mother of all BCI
Benchmarks \cite{moabb} (MOABB) platform. A previous intern from 2024, Racim
Menasria, had to reproduce DreamDiffusion. He choosed to use the
GrosseWentrup2009 \cite{GrosseWentrup2009} and Schirrmeister2017
\cite{Schirrmeister2017} datasets. In the process of reproducing what Racim did,
an error was corrected in a dataset with a GitHub pull request
\url{https://github.com/NeuroTechX/moabb/pull/781}. Not only the datasets used
were hard to find but they were also very long to download as the server
delivering them had internet limitations.

Many little problems made the code work on the authors' machine but not on any
other machine, those little issues were fixed and a pull request was proposed,
but the authors did not answer
\url{https://github.com/bbaaii/DreamDiffusion/pull/42}.

At one time the goal was to try and improve DreamDiffusion with novels idea and
it was needed to go through the whole code to begin with it. The code is taken
from \cite{mindvis} so and the structure of the files and their roles was new so
it also took a good amount of time to understand every bit.

\subsubsection{Results}

\insertfigure
    {assets/images/dd-good}
    {width=0.7\textwidth}
    {Good result, ground truth on the left}
    {fig:dd-good}

It happens quite often that the generated image is of the good class (around 30
percent of the time), but it fails to show good similarities, it feels like it's
an EEG classifier and then an image generator from the guessed class. In
\cref{fig:dd-good}, there is not many similarities apart from the fact that
these are clownfish pictures, this is probably due to a not high enough
coefficient to take into account the diffusion loss. This also explains why some
images are incoherent.

We can distinguish two common mistakes that happen in the 70 percents left. The
first one is when the generated image is coherent and looks like a real image,
but from another class \cref{fig:dd-coherent}. The other mistake is when the
generated image is not coherent, it's hard to understand what is shown
\cref{fig:dd-incoherent}.

\insertfigure
    {assets/images/dd-coherent}
    {width=0.7\textwidth}
    {Coherent generated image, bad class (ground truth on the left)}
    {fig:dd-coherent}

\insertfigure
    {assets/images/dd-incoherent}
    {width=0.7\textwidth}
    {Incoherent generated image (ground truth on the left)}
    {fig:dd-incoherent}